\chapter{Déconvolution d'un signal triangulaire}

Le problème considéré est celui de la reconstruction d'un signal $x$, supposé inconnu car non accessible à une mesure directe, ayant subi un filtrage gaussien inhérent à l'instrument de mesure. Il s'agit d'un filtrage linéaire d'usage courant en traitement d'image; il est utilisé, ici, dans un contexte monodimensionnel. La réponse impulsionnelle (RI) d'un filtre gaussien est définie par:

\eql{
	\label{eq:filtre}
	h(t) = \frac{1}{\sigma_h \sqrt{2\pi}} \exp\left(-\frac{(t-m_h)^2}{2\sig_h^2}\right)
}

Le signal $y$ mesuré en sortie du filtre, seul accessible, est supposé être corrompu par des incertitudes additives. En effet, les mesures en sortie de filtre sont à minima sujettes au bruit de
quantification dû au convertisseur analogique-numérique de la chaîne de mesure.

Il s'agit d'appréhender le caractère instable de l'opération de déconvolution, lorsque celle-ci est effectuée de façon  \guillemotleft naïve\guillemotright, soit par division polynomiale (fonction \texttt{deconv} de Matlab), soit par filtrage inverse (dans le domaine fréquentiel ou par moindres carrés qui est l'équivalent matriciel). En présence d'erreurs (de bruit) sur les données, même très faibles, les résultats fournis par ces
méthodes d'inversion peuvent, en effet, être inexploitables. On parle problème inverse \emph{mal posé}. La réponse en fréquence du système à inverser permet de caractériser le degré de difficulté du problème. L'exercice illustrera la (les) solution(s) obtenues par régularisation linéaire quadratique.

\section{Construction du jeu de données}

La fréquence d'échantillonnage est fixée à la valeur $fe = 1$ Hz. Les échantillons du signal d'entrée sont notés $x(n)$, ceux de la réponse impulsionnelle du filtre $h(n)$ et ceux du signal mesuré en sortie $y(n)$. Le signal mesuré en sortie est sujet à des incertitudes additives notées $b(n)$ et le problème direct peut alors être modélisé comme suit :
\eql{
	\label{eq:modele}
	y(n) = (x \ast h)(n) + b(n).
}

\paragraph{Simulation du signal d'entrée.} Le signal d'entrée, supposé être triangulaire, est défini par
\eq{
	x(t) = \left\{
		\begin{aligned}
			&\frac{t}{L_x/2}, \quad \forall t \in [0,\frac{L_x}{2}]\\
			&\frac{-t+L_x}{L_x/2}, \quad \forall t \in [\frac{L_x}{2}, L_x].
		\end{aligned}
		\right.
}  
Simuler le signal à temps discret $x(n)$ correspondant, pour une durée d'observation $L_x = 100$s.

\paragraph{Aide:} remplacer $t$ par $nT_e$ dans l'équation définissant $x(t)$. Il est possible d'utiliser une instruction du type \texttt{x = [x1,x2]}, où \texttt{x1} et \texttt{x2} représentent deux tronçons successifs du signal. De manière générale, il est conseillé d'apdapter le script ci-dessous qui permet de maîtriser les différents paramètres de la discrétisation.

\begin{lstlisting}
Lx = 100;	% duree d'observation du signal d'entree (secondes)
fe = 2;		% frequence d'echantillonnage (Herz)
Te = 1/fe;	% periode d'echantillonnage (secondes)
Nx = Lx/Te;	% nombre d'échantillons du signal
nx = 0:Nx-1;	% indice temporel (entier)
tx = nx*Te;		% base de temps (secondes)
\end{lstlisting}

\paragraph{Simulation du filtre.} Pour les besoins de la simulation, déterminer la réponse impulsionnelle $h(n)$ du filtre gaussien obtenue par discrétisation de l'équation \eqref{eq:filtre}. Le filtre est de plus supposé causal et ses paramètres seront pris égaux respectivement à $m_h = 15$ et $\sig_h = 4$. Enfin, la réponse impulsionnelle est tronquée à un horizon de 30 s. Représ˚enter cette réponse impulsionnelle.

\medbreak

\noindent \textbf{Rappel:} la réponse impulsionnelle d'un filtre causal est nul aux temps négatifs


\paragraph{Simulation du problème direct, première approche: fonction \texttt{conv}.} Comme le filtre gaussien est un filtre \textbf{linéaire invariant}, il est possible de calculer sa sortie par le produit de convolution discret donné par l'équation \eqref{eq:modele}.
\begin{enumerate}
	\item Simuler le signal de sortie non bruité $\bm{y}_{nb}(n)$ à l'aide de la fonction \texttt{conv}.
	\item Si le signal de sortie $\bm{y}$ est mesuré à l'aide d'un convertisseur analogique-numérique, il est sujet au bruit de quantification. L'opération de quantification peut être simulée par l'instruction \texttt{y=round(ynb*2\^{}10)/2\^{}10}. Simuler le signal quantifié y et représenter celui-ci à l'aide de la fonction \texttt{stairs}.
	\item Déterminer et représenter le bruit de quantification $b(n)$.
\end{enumerate}

\paragraph{Simulation du problème direct, deuxième approche: formulation matricielle.} Dans cette partie, le signal de sortie non bruité est simulé par l'équation matricielle $\bm{y}_{nb} = \bm{H}\bm{x}$. L'usage d'un paramétrage soigné est recommandé pour permettre de modifier aisément les conditions de simulation et d'affichage des résultats.
\begin{enumerate}
	\item Déterminer la dimension Ny du vecteur de sortie ynb à partir des dimensions $N_x$ et $N_h$ des vecteurs d'entrée $\bm{x}$ et de réponse impulsionnelle $\bm{h}$ afin d'obtenir le même nombre d'échantillons que par la fonction \texttt{conv}.
	\item Construire la matrice de convolution \bm{H} à l'aide de la fonction \texttt{toeplitz}. Pour cela, il est proposé d'initialiser le premier vecteur ligne et le premier vecteur colonne de H avec des vecteurs nuls de dimension appropriée à l'aide la fonction \texttt{zeros}, puis de remplacer certains de leurs éléments avec les éléments de \bm{h}.
	\item Déterminer la sortie non bruitée par une opération matricielle, puis appliquer la fonction permettant de simuler l'erreur de troncature.
\end{enumerate}


\section{Déconvolution}

\subsection{Déconvolution \guillemotleft naïve\guillemotright et constat de l'instabibilité}
\begin{enumerate}
	\item Reconstruire le signal d'entrée en utilisant l'opération de déconvolution Matlab (\texttt{deconv}), en se plaçant, dans un premier temps, dans le cas idéal où l'on a accès à la sortie non bruitée $\bm{y}_{nb}$.
	\item À présent, reconstruire le signal d'entrée à partir des échantillons mesurés du signal de sortie \bm{y} (cadre réaliste). Le signal reconstruit sera appelé $\bm{x}_{rec}$. Que vaut le RSB du signal de sortie?
	\item Étudier de façon empirique l'influence de la \guillemotleft forme\guillemotright de la réponse impulsionnelle sur la qualité de la reconstruction en modifiant la valeur du paramètre $\sigma_h$.
\end{enumerate}

\subsection{Interprétation fréquentielle de l'inversion}

\subsection{Déconvolution par moindres carrés ($\ell_2$)}
\begin{enumerate}
	\item Calculer le signal déconvolué à l'aide de l'estimateur des moindres carrés.
	\item Tracer le signal déconvolué et calculer l'énergie de l'erreur de reconstruction.
	\item Tracer le spectre des valeurs singulières de la matrice $\bm{H}^\top \bm{H}$ et déterminerson conditionnement.
	\item Comment le conditionnement évolue-t-il en fonction de la valeur du paramètre $\sigma_j$ du filtre?
\end{enumerate}

\subsection{Déconvolution par régularisation linéaire quadratique ($\ell_2\ell_2$)}
Cette partie concerne la mise en ÷uvre de la solution obtenue par minimisation d'un critère composite de fonctionnelles $L$ et $R$ quadratiques : $J(\bm{x}) = L(\bm{x}) + \lambda R(\bm{x})$.
\begin{enumerate}
	\item Construire la matrice carrée $\bm{D}_1$, de dimension $N_x$, telle que le produit $\bm{D}_1x$ donne les
différences premières\footnote{Les différences premières, lorsqu'elles sont divisées par les pas $T_e$, correspondent à une approximation de la dérivée première.}: $x(n + 1) − x(n)$.
\medbreak
\textbf{Aide:} On s'aidera du code suivant
\begin{lstlisting}
dcol1 = zeros(1:Nx);
dcol1(1:2) = [1 -1];
dlig1 = zeros(1,Nx);
dlig1(1,1) = dcol(1,1);
D1 = toeplitz(dcol1,dlig1);
\end{lstlisting}
\medbreak
\textbf{Variante:} remplacer la seconde ligne du script ci-dessus par
\begin{lstlisting}
dcol1(1:3) = [1 -2 1];
\end{lstlisting}
donne accés aux différences secondes définissant alors la matrice $\bm{D}_2$. 
	\item Calculer la solution régularisée en déterminant le coefficient de régularisation $\lambda$ de manière empirique\footnote{Par exemple, en testant plusieurs valeurs dans un intervalle conséquent, comme $[10^{-7},10^2]$.}.
	\item En simulation, le signal à reconstruire est parfaitement connu, il est donc possible de déterminer la valeur du coefficient de régularisation qui minimise l'énergie de l'erreur de
reconstruction. Tracer l'évolution de l'énergie de l'erreur de reconstruction en fonction de $\lambda$. Ecrire une procédure permettant d'approcher la valeur optimale de $\lambda$, par recherche exhaustive sur un intervalle, à l'aide de la fonction \texttt{min}. 
	\item Tracer la \guillemotleft courbe en L\guillemotright en représentant, sur une échelle log-log (fonction \texttt{loglog}), la fonctionnelle de régularisation $R(\hat{\bm{x}}(\la) = \norm{\bm{D}_1\hat{\bm{x}}(\la)}_2^2$ en fonction du critère d'adéquation aux données $L(\hat{\bm{x}}(\la)) = \norm{\bm{y} - \bm{H}\hat{\bm{x}}(\la)}_2^2$ (moindres carrés).
	\medbreak
	\textbf{Aide:} vous pouvez faire varier $\lambda$ de la manière suivante
\begin{lstlisting}
expmin = -7;
expmax = 2;
pas = 0.1;
id = 0;
for p = expmin:pas:expmax
	lambda = 10^p;
	id = id+1;
	...
\end{lstlisting}
	\item Que se passe-t-il lorsque l'on remplace la matrice $\bm{D}_1$ par la matrice identité $I_{N_x}$ ou par $\bm{D}_2$? Quelle pénalisation constitue une façon appropriée de coder les connaissances a priori dont on dispose sur le signal à reconstruire.
\end{enumerate}

